% !TEX root = ../../../main.tex
% 所属：第四部分：Fourier 分析 / 第2章：Fourier级数的收敛性
% 节：2.3 一个反例——连续函数的 Fourier 级数可能发散
% 正文按原文件顺序保存；由 ../chapter2.tex 调用。
\section{一个反例——连续函数的 Fourier 级数可能发散}


上一节说明,若函数在一点可微,则其 Fourier 级数在该点收敛到函数值.下面构造一个连续周期函数,使其 Fourier 级数在指定点发散.这表明,上一节中的可微条件不能直接削弱为连续条件.

\begin{proposition}
	存在圆周上的连续函数 $F$,使得其 Fourier 部分和在原点无界.更准确地说,存在严格递增的正整数列 $\{M_k\}$,满足
	\[
	\abs{S_{M_k}(F)(0)}\longrightarrow\infty.
	\]
	因此,$F$ 的 Fourier 级数在 $\theta=0$ 处发散.通过平移,也可以使发散发生在任意预先指定的点 $\theta_0$.
\end{proposition}

构造的关键是 Fourier 部分和中的\textbf{对称性破缺}.通常
\[
S_N(f)(\theta)=\sum_{-N}^{N}\widehat f(n)e^{in\theta}
\]
同时截取正频率与负频率,二者之间可能产生显著抵消.我们将先构造一个一致有界的对称频率块,再通过适当移动频率,使某个 Fourier 部分和恰好从中间切开该频率块.被截出的单边部分将在原点具有对数级增长.



考虑 $h$ ,圆周上的锯齿函数：它关于 $\theta$ 为奇函数,并且当 $0<\theta<\pi$ 时,
\[
h(\theta)=i(\pi-\theta).
\]
由前面的计算可知,
\begin{equation}\label{eq:sawtooth-fourier}
	h(\theta)\sim\sum_{n\neq0}\frac{e^{in\theta}}{n}.
\end{equation}
式 \eqref{eq:sawtooth-fourier} 中正、负频率关于原点对称.若只保留负频率,得到形式级数
\[
\sum_{n=-\infty}^{-1}\frac{e^{in\theta}}{n}.
\]
这个单边级数不可能是某个 Riemann 可积函数的 Fourier 级数.事实上,若它是某个函数 $\widetilde h$ 的 Fourier 级数,则 $\widetilde h$ 有界,且在原点的 Abel 平均满足
\[
\abs{A_r(\widetilde h)(0)}
=
\sum_{n=1}^{\infty}\frac{r^n}{n}
\longrightarrow\infty,
\qquad r\to1^-.
\]
另一方面,由 Abel 平均的卷积表示及 Poisson 核的正性,
\[
\abs{A_r(\widetilde h)(0)}
=
\abs{(\widetilde h\ast P_r)(0)}
\leq
\sup_{\theta}\abs{\widetilde h(\theta)},
\]
得到矛盾.

对每个正整数 $N$,定义三角多项式

	\[ f_N(\theta)=\sum_{1\leq\abs{n}\leq N}\frac{e^{in\theta}}{n},
	\qquad
	\widetilde f_N(\theta)=\sum_{-N\leq n\leq-1}\frac{e^{in\theta}}{n}. \]

在原点,
\[
\abs{\widetilde f_N(0)}
=
\sum_{n=1}^{N}\frac1n
\geq
\int_1^N\frac{\dd x}{x}
=
\log N.
\]
因此存在绝对常数 $c>0$,使得
\begin{equation}\label{eq:one-sided-log-growth}
	\abs{\widetilde f_N(0)}\geq c\log N.
\end{equation}
与此相反,\textbf{完整的对称和 $f_N$ 关于 $N$ 与 $\theta$ 一致有界}.为证明这一点,先给出一个简单引理.

\begin{lemma}
	设复数列 $\{c_n\}_{n\geq1}$ 的 Abel 平均
	\[
	A_r=\sum_{n=1}^{\infty}r^nc_n,
	\qquad 0<r<1,
	\]
	在 $r\to1^-$ 时有界.若 $c_n=O(1/n)$,则其部分和
	\[
	T_N=\sum_{n=1}^{N}c_n
	\]
	也是有界的.
\end{lemma}

\begin{proof}
	取 $r=1-1/N$.选择常数 $M>0$,使得对所有 $n$ 及充分接近 $1$ 的 $r$ 都有
	\[
	n\abs{c_n}\leq M,
	\qquad
	\abs{A_r}\leq M.
	\]
	于是
	\[
	S_N-A_r
	=
	\sum_{n=1}^{N}(1-r^n)c_n
	-
	\sum_{n=N+1}^{\infty}r^nc_n.
	\]
	利用
	\[
	1-r^n=(1-r)(1+r+\cdots+r^{n-1})\leq n(1-r),
	\]
	得到
	\[
	\begin{aligned}
		\abs{S_N-A_r}
		&\leq
		\sum_{n=1}^{N}\abs{c_n}(1-r^n)
		+
		\sum_{n=N+1}^{\infty}r^n\abs{c_n}\\
		&\leq
		M\sum_{n=1}^{N}(1-r)
		+
		\frac{M}{N}\sum_{n=N+1}^{\infty}r^n\\
		&\leq
		MN(1-r)+\frac{M}{N(1-r)}
		=2M.
	\end{aligned}
	\]
	故 $\abs{S_N}\leq3M$,从而部分和有界.
\end{proof}

将引理应用于
\[c_n=\frac{e^{in\theta}}{n}+\frac{e^{i(-n)\theta}}{-n}\]
则
\[
f_N(\theta)=\sum_{n=1}^{N}c_n(\theta).
\]
显然 $n\abs{c_n(\theta)}\leq2$.另一方面,相应的 Abel 平均正是锯齿函数 $h$ 的 Abel 平均,
\[
\sum_{n=1}^{\infty}r^nc_n(\theta)
=A_r(h)(\theta)
=(h\ast P_r)(\theta),
\]
故
\[
\abs{A_r(h)(\theta)}\leq\sup_{\varphi}\abs{h(\varphi)}.
\]
上述估计与 $r,\theta$ 无关.因此存在绝对常数 $C>0$,使得
\begin{equation}\label{eq:fN-uniform-bound}
	\abs{f_N(\theta)}\leq C,
	\qquad N\geq1,\quad\theta\in\R.
\end{equation}



定义
\begin{equation}\label{eq:shifted-blocks}
	P_N(\theta)=e^{i2N\theta}f_N(\theta),
	\qquad
	\widetilde P_N(\theta)=e^{i2N\theta}\widetilde f_N(\theta).
\end{equation}
乘以 $e^{i2N\theta}$ 会把所有频率向右移动 $2N$ 个单位.因此,$P_N$ 的非零 Fourier 系数只出现在
\[
N\leq n\leq3N,
\qquad n\neq2N.
\]
其中区间 $[N,2N-1]$ 对应原来的负频率,区间 $[2N+1,3N]$ 对应原来的正频率.

\begin{lemma}\label{lem:frequency-block-cut}
	对任意正整数 $M,N$,有
	\[
	S_M(P_N)=
	\begin{cases}
		P_N,&M\geq3N,\\
		\widetilde P_N,&M=2N,\\
		0,&M<N.
	\end{cases}
	\]
\end{lemma}

\begin{proof}
	当 $M<N$ 时,截断区间 $[-M,M]$ 不包含 $P_N$ 的任何频率,故 $S_M(P_N)=0$.当 $M\geq3N$ 时,全部频率都被保留,故 $S_M(P_N)=P_N$.最后,当 $M=2N$ 时,仅保留区间 $[N,2N-1]$ 中的频率,这恰好是 $\widetilde f_N$ 的频率向右移动 $2N$ 后所得的 $\widetilde P_N$.
\end{proof}

注意,由式 \eqref{eq:fN-uniform-bound},
\begin{equation}\label{eq:PN-uniform-bound}
	\abs{P_N(\theta)}=\abs{f_N(\theta)}\leq C,
\end{equation}
但由式 \eqref{eq:one-sided-log-growth},
\begin{equation}\label{eq:Ptilde-log-growth}
	\abs{\widetilde P_N(0)}
	=
	\abs{\widetilde f_N(0)}
	\geq c\log N.
\end{equation}
因此,$P_N$ 本身一致有界,而部分和算子 $S_{2N}$ 从频率块中截出的部分却在原点具有对数级增长.



\begin{proof}[命题的证明]
	选择正数列与正整数列
	\[
	\alpha_k=\frac1{k^2},
	\qquad
	N_k=3^{2^k}.
	\]
	则
	\[
	\sum_{k=1}^{\infty}\alpha_k<\infty,
	\qquad
	N_{k+1}>3N_k,
	\qquad
	\alpha_k\log N_k
	=
	\frac{2^k\log3}{k^2}
	\longrightarrow\infty.
	\]
	定义
	\begin{equation}\label{eq:continuous-divergent-example}
		F(\theta)=\sum_{k=1}^{\infty}\alpha_kP_{N_k}(\theta).
	\end{equation}
	由式 \eqref{eq:PN-uniform-bound} 与 Weierstrass 判别法,级数 \eqref{eq:continuous-divergent-example} 一致收敛.因此 $F$ 是圆周上的连续函数.
	
	由于 $N_{k+1}>3N_k$,不同 $P_{N_k}$ 的频率支集彼此分离.固定 $m$,考察阶数为 $2N_m$ 的 Fourier 部分和.由引理 \ref{lem:frequency-block-cut},当 $k>m$ 时,
	\[
	S_{2N_m}(P_{N_k})=0;
	\]
	当 $k<m$ 时,
	\[
	S_{2N_m}(P_{N_k})=P_{N_k};
	\]
	而当 $k=m$ 时,
	\[
	S_{2N_m}(P_{N_m})=\widetilde P_{N_m}.
	\]
	又因为
	\[
	P_{N_k}(0)=f_{N_k}(0)
	=
	\sum_{1\leq\abs{n}\leq N_k}\frac1n
	=0,
	\]
	所以
	\[
	S_{2N_m}(F)(0)
	=
	\alpha_m\widetilde P_{N_m}(0).
	\]
	由式 \eqref{eq:Ptilde-log-growth},
	\[
	\abs{S_{2N_m}(F)(0)}
	\geq
	c\alpha_m\log N_m
	=
	c\frac{2^m\log3}{m^2}
	\longrightarrow\infty.
	\]
	故 $F$ 的 Fourier 部分和在原点无界,从而其 Fourier 级数在原点发散.
	
	若希望发散发生在预先指定的点 $\theta_0$,只需令
	\[
	F_{\theta_0}(\theta)=F(\theta-\theta_0).
	\]
	此时
	\[
	S_N(F_{\theta_0})(\theta_0)=S_N(F)(0),
	\]
	结论成立.
\end{proof}
